\documentclass[letterpaper,11pt]{article}
\usepackage[top=0.8in, bottom=0.8in, left=0.9in, right=0.9in]{geometry}
\usepackage[hidelinks]{hyperref}
\usepackage{lmodern}
\usepackage{parskip}
\usepackage{microtype}
\setlength{\parskip}{6pt}
\setlength{\parindent}{0pt}
\begin{document}
\begin{flushleft}
\textbf{\large Ajay Venkatesh} \\
Brooklyn, NY \\
+1 516 585 9013 \\
\href{mailto:av3855@nyu.edu}{av3855@nyu.edu}
\end{flushleft}
\vspace{10pt}
May 21, 2026
\vspace{6pt}

Hiring Manager \\
iCapital \\
New York, NY

\vspace{10pt}

Dear Hiring Manager,

My name is Ajay Venkatesh. I'm finishing my M.S. in Computer Engineering at NYU, with production software experience at PwC in Bengaluru, a summer SDE internship at Amazon, and ongoing on-campus work at NYU's Novel AI Technologies. I'm writing about the Associate Full Stack Engineer role on the Post-Investments Technology team.

The strongest match is the frontend-heavy full-stack work I've shipped. At Amazon I built a \textbf{React 18} frontend in \textbf{TypeScript} with optimized API integration, advanced filtering, and client-side caching against a Java microservice backend, sitting behind a careful auth model and role-based access. At NYU I shipped a customer-facing insurance analytics platform on \textbf{React, Node.js, TypeScript, REST APIs}, backed by \textbf{PostgreSQL}, deployed via \textbf{Docker} on \textbf{EC2}, with \textbf{RBAC}, row-level security, and Stripe billing. The work was document-heavy and workflow-driven, since insurance underwriting is exactly that kind of UI, and I designed the data model end-to-end as the product evolved.

At Campus Mesh I extended this into \textbf{Node.js microservices} powering real-time academic workflows, with \textbf{MySQL (Sequelize)} and \textbf{DynamoDB} stores, a serverless backend on AWS \textbf{Lambda} and \textbf{API Gateway}, and partitioned data modeling for high-throughput chat. That gives me the muscle for both sides of a microservices-based architecture, which is the direction the role grows into over time.

On overlap with your stack: I'm fluent in \textbf{TypeScript}, modern frontend frameworks (\textbf{React}, \textbf{Next.js}), \textbf{REST API design}, \textbf{microservices architecture}, and \textbf{AWS} (\textbf{S3}, \textbf{Lambda}, \textbf{API Gateway}, \textbf{IAM}). I work through code review as part of daily practice and I'm comfortable across the stack from data layer to UI. \textbf{Scala} and \textbf{Ruby on Rails} are the two genuine gaps, and I'd plan to come up the curve on them quickly inside the existing backend codebase while contributing in earnest on the TypeScript frontend from day one. I've picked up new backend stacks before, including Java at Amazon and Node.js at Campus Mesh, and I'm comfortable doing it again.

What pulls me toward iCapital is the surface. Alternative investments, structured products, and annuities are some of the harder financial workflows to make legible to humans, and that's exactly where a careful frontend earns its keep. I'd be excited to build client-facing platforms where a small UX improvement can compound across thousands of advisors and very large amounts of capital, and to do that work here in New York alongside a team that takes the engineering side seriously.

I'd welcome the chance to discuss further. Thanks for considering my application.

\vspace{12pt}
Sincerely, \\
Ajay Venkatesh

\end{document}
